% !TeX root = ../papers/nc-p3.tex
\section{Euler導分による中心の持ち上げと十五次元成分の正則開集合}
\label{sec:quartic-euler-lifting}
$B$を四頂点で折り畳んだ経路長次数付き代数とし、$E(x)=n x$ ($x\in B_n$)を
Euler導分とする。以下の導分は各頂点冪等元を零に送る。従って内部次数$d$の
内導分の係数は$\bigoplus_i e_i B_d e_i$から取る。

\begin{proposition}\label{prop:euler-central-lifting}
体の標数を零、$d>0$とし、$B$の内部次数$d$の全導分が内導分と
$Z(B)_d E$の和であるとする。任意の局所完備Noether環上の次数別平坦な
斉次変形で、中心の次数$d$の非零因子は、同次数の中心元へ持ち上がる。
\end{proposition}
\begin{proof}
極大イデアルを$\mathfrak m$とし、$z\in Z(B)_d$の中心リフトが
$\mathfrak m^n$を法として得られたとする。その任意の次のリフト$y$を
頂点対角に選ぶ。$[y,-]$は$I=\mathfrak m^n/\mathfrak m^{n+1}$に値を持つ。
次数別平坦性と下の次数の中心性により、これはリフトの選択に関して
入力側で剰余体に降り、$B\to I\otimes B$の導分となる。
各$\ell\in I^*$で押し出した導分を$[v,-]+hE$と書く。
特に$z$のリフトとして$y$自身を使えば値は零である。
従って$0=dhz$であり、非零因子性から$h=0$である。
$I$の基底を取り、得られた内導分係数を$y$から引けば
$\mathfrak m^{n+1}$を法としても中心となる。完備化上では逆極限を取る。
ここで変形は頂点分解と次数を保ち、各斉次成分は有限階数とする。
\end{proof}

\begin{theorem}\label{thm:full-component-central-open}
十五次元成分の反対階数16の有理点$t$を含む一つの非空ザリスキー開集合で、
\emph{全ての幾何点}が標準Koszul左右AS正則となる。四次中心は五次元であり、
四つの中心四次式が生成する多項式環上で各始点が階数256の自由加群である。
\end{theorem}
\begin{proof}
$t$では既にこの自由表示と四次中心の文字次数別次元$(3,2,0,0)$を知っている。
次数4の導分は各矢を同じ端点の次数5成分に送り、六関係にはLeibniz則を課す。
文字次数$h$を固定すると各始点の四矢で合計56係数、四始点で224未知数である。
各関係の二文字の片方を順に五文字の商基底で置き換え、六文字商へ射影する。
\path{computations/verify_quartic_euler_derivations.py}が有理経路基底を
素数1009で還元し、$h=0,1,2,3$の順に階数$188,184,192,192$を検算する。
よって導分空間の次元の上界は$(36,40,32,32)$である。
他方、四次内導分の次元は$(33,38,32,32)$であり、各文字の四次中心の
Euler倍はこれと独立である。実際、中心四変数の一つに$hE$を作用させれば、
その非零因子と$4h$の積を得るが、内導分は中心を消す。各上界は従って等号である。

テンソル成分の$t$での完備局所環を取る。各Artin商上で隣接二切断の六次主小行列は
単元のままで、節\ref{sec:unrestricted-doublet}の次数別Koszul複体の分裂により
経路代数は次数別に自由である。命題\ref{prop:euler-central-lifting}を適用できる。
四次中心の基底は全て非零因子に取れる。最初の四つは認証した多項式変数とし、
残る元$z$にはその一変数$x$の一般定数倍を足す。多項式環の分数体上、
$x$の乗法は可逆なので$z+c x$の乗法の行列式は$c$について零多項式ではない。
元の代数も多項式環上で自由なため、適切な$c$で非零因子となる。
従って五つの独立な中心元が完備局所環上で同時に延びる。

次数4・5の経路空間を自由化した近傍で中心方程式はベクトル束の写像である。
完備化で五つの核切断は剰余でも独立なので、核次元が5未満となるために必要な
小行列式は完備環で全て零である。局所環から完備化への忠実平坦性により
それらは局所環でも零となる。一方$t$での階数を与える小行列式は単元である。
有限個の係数しか使わないため開近傍で四次中心は階数5の部分束となる。

最初の四変数をこの束の切断として取り、割った商の次数13までの上界が
$(1+t+t^2+t^3)^4$の係数となる有限個の開条件を課す。
$t$で全てが成立する。元の経路商の各次数には、この既約成分の正則な
幾何学的一般点からの標準Hilbert関数の下界がある。四中心式の多項式環からの
自由加群による全射とこの下界を比べ、定理\ref{thm:klein-rational-period-four}の
左右完全性の議論を適用すれば、同じ一つの開集合の全幾何点で結論を得る。
最後に開集合の基底変換による移動を合わせることで、印なしの主張となる。
\end{proof}

\begin{corollary}
この成分の閉包の標準AS点は全て多項式恒等式を持つ。
\end{corollary}
\begin{proof}
各終点の列は開集合で中心四変数上自由階数256である。
系\ref{cor:klein-pi-specializations}と同様に標準恒等式$s_{512}=0$を、
斉次な語に代入し、必要な有限個の経路次数の定階数性によって各標準AS点へ特殊化する。
\end{proof}
